\chapter{Conclusiones}

El trabajo experimental realizado en el laboratorio permitió verificar en forma directa el funcionamiento del módulo
matemático FFT del osciloscopio digital, contrastar las ventajas del análisis espectral frente a la visualización
temporal y responder a los aspectos centrales planteados en la práctica:

\section{Ventajas y Desventajas del Análisis en Frecuencia frente al Tiempo}

\noindent\textbf{Ventajas del Dominio de la Frecuencia:}
\begin{itemize}[noitemsep, topsep=2pt]
    \item \textbf{Gran rango dinámico y discriminación de armónicos:} Permite identificar componentes débiles y
    productos espurios situados a $\qty{-40}{\decibel}$ o $\qty{-50}{\decibel}$ por debajo de la fundamental, que
    en el tiempo quedan enmascarados dentro del trazo.
    \item \textbf{Medición cuantitativa de distorsión y modulación:} Facilita el cálculo directo de la distorsión
    armónica total (THD), del rechazo a intermodulación (IMD3) y de los índices de modulación en AM y FM.
    \item \textbf{Deducción de parámetros temporales:} A partir del espaciamiento entre rayas armónicas se calcula el
    período ($T = 1/\Delta f_n$) y a partir de los nulos de la envolvente sinc se obtiene el ancho de pulso
    ($t_0 = 1/\Delta f_{\text{min}}$).
\end{itemize}

\noindent\textbf{Desventajas del Dominio de la Frecuencia:}
\begin{itemize}[noitemsep, topsep=2pt]
    \item \textbf{Pérdida de información de fase:} La FFT exhibe magnitudes espectrales, perdiendo la fase relativa
    entre armónicos e impidiendo reconstruir la forma de onda instantánea.
    \item \textbf{Compromiso de resolución y fuga espectral:} La resolución depende del tiempo de registro ($\Delta f = 1/T_{\text{rec}}$).
    Si la señal no es periódica en la ventana de adquisición, se produce fuga espectral (\emph{leakage}).
    \item \textbf{Vulnerabilidad a frecuencias alias:} Si la tasa de muestreo no respeta el criterio de Nyquist,
    surgen componentes falsas replegadas que inducen a errores de diagnóstico.
\end{itemize}

\section{Ventajas y Desventajas de los Distintos Tipos de Ventanas}
Cada ventana matemática de ponderación establece un compromiso entre resolución en frecuencia y exactitud en amplitud,
resumido en el Cuadro~\ref{tab.concl:ventanas}:

\begin{table}[!ht]
    \centering
    \resizebox{\textwidth}{!}{%
    \begin{tabular}{lccc}
        \toprule
        \textbf{Ventana} & \textbf{Resolución de Frecuencia} & \textbf{Fuga Espectral (Lóbulos)} & \textbf{Aplicación Recomendada} \\
        \midrule
        Rectangular & Máxima (lóbulo muy estrecho) & Alta (lóbulos a $\qty{-13}{\decibel}$) & Señales transitorias y pulsos \\
        Hanning & Media (buen compromiso) & Baja (lóbulos a $\qty{-32}{\decibel}$) & Señales periódicas y ruido \\
        Flattop & Menor (lóbulo ancho) & Mínima (cima plana, $< \qty{0.1}{\decibel}$) & Medición exacta de amplitudes y THD \\
        \bottomrule
    \end{tabular}}
    \caption{Comparación de ventajas y aplicaciones de las ventanas Rectangular, Hanning y Flattop.}
    \label{tab.concl:ventanas}
\end{table}

\section{Análisis del Fenómeno de Alias (Experimento 3)}
Al fijar $f_s = \qty{25.0}{\kilo\Sa\per\second}$ ($f_{\text{Nyq}} = \qty{12.5}{\kilo\hertz}$), se constató que:
\begin{itemize}[noitemsep, topsep=2pt]
    \item Para $f_{\text{in}} < \qty{12.5}{\kilo\hertz}$, la presentación espectral fue unívoca y correcta.
    \item Al inyectar una senoidal de $\qty{15.20}{\kilo\hertz}$, la componente se replegó especularmente sobre la banda base en $f_{\text{alias}} = |f_s - f_{\text{in}}| = \qty{9.80}{\kilo\hertz}$, evidenciando la aparición de una raya falsa sin aviso del instrumento.
    \item Con onda cuadrada, las múltiples componentes que excedieron Nyquist colapsaron el espectro con un peine de rayas espurias, ratificando la necesidad de filtros antialiasing o sobremuestreo adecuado.
\end{itemize}

\section{Diferencias entre los Modos de Adquisición}
\begin{itemize}[noitemsep, topsep=2pt]
    \item \textbf{Normal:} Muestreo convencional; refleja la variabilidad del piso de ruido en tiempo real.
    \item \textbf{Pico (Peak Detect):} Retiene máximos y mínimos analógicos; captura transitorios rápidos a costa de elevar el piso de ruido.
    \item \textbf{Promedio (Average):} Promedia sincrónicamente múltiples adquisiciones (64 cuentas), reduciendo la varianza estocástica del ruido y estabilizando trazas de baja amplitud.
\end{itemize}

\section{Funcionamiento del Modulador a Diodo (Experimentos 4 y 5)}
A partir de la alinealidad del diodo ($i_D(v) \approx a_1 v + a_2 v^2 + a_3 v^3 + \dots$):
\begin{itemize}[noitemsep, topsep=2pt]
    \item El término cuadrático ($a_2 v^2$) produce la modulación AM ($f_c \pm f_m$), seleccionada por el tanque sintonizado en $\approx \qty{50}{\kilo\hertz}$.
    \item El término cúbico ($a_3 v^3$) engendra productos de intermodulación ($2f_1 - f_2$ y $2f_2 - f_1$), verificados en el Experimento~5 con un rechazo de $\qty{35.2}{\decibel}$ respecto a los tonos fundamentales.
\end{itemize}

\section{Conclusión General}
El desarrollo del trabajo práctico permitió corroborar experimentalmente el análisis espectral mediante FFT, contrastar el comportamiento de las ventanas matemáticas, evidenciar los efectos del submuestreo y cuantificar la acción linealizante de la realimentación negativa sobre la distorsión armónica en amplificadores clase A, con estricto seguimiento de las pautas metrológicas de la cátedra.
